% GENERATED, edit the French unit or the translation table
% GENERATED from content/data/enrichment.json by mxbuild.enrichment; do not edit.
\mxunit{enrichment-diophantine-notes}{Bézout and integer equations}
\mxsubtitle{Advanced high school · extension \textperiodcentered{} Lesson and method \textperiodcentered{} 04}

\begin{prerequisites}
Euclidean algorithm and divisibility in $\Z $.
\end{prerequisites}

\begin{keypoint}[Ideas and method]
For positive integers a and b, the Euclidean algorithm computes $d = \pgcd(a,b)$. Back-substitution expresses d as au + bv for integers u and v: a Bézout identity.

The equation ax + by = c has integer solutions exactly when d divides c. Necessity follows because d divides ax + by; sufficiency follows by multiplying a Bézout identity by $\frac{c}{d}$.

If $\pt{x_{0};y_{0}}$ is one solution, all solutions are $x = x_{0} + (\frac{b}{d})t$ and $y = y_{0} - (\frac{a}{d})t$ for integer t. Conditions $x \geq  0$ and $y \geq  0$ must then be imposed separately.
\end{keypoint}

\begin{example}
For $15x + 21y = 6$, use $21 = 15 + 6$ and $15 = 2 \times  6 + 3$. Thus $3 = 3 \times  15 - 2 \times  21$. One solution is $\pt{6;-4}$, and all solutions are $x = 6 + 7t, y = -4 - 5t$.
\end{example}

\begin{activity}[Explore]
Try several values of t in the example and verify the equation exactly.
\end{activity}

